\documentclass[10pt,twocolumn,a4paper]{article}

\usepackage[margin=0.75in]{geometry}
\usepackage{times}
\usepackage{enumitem}
\usepackage{amsmath}
\usepackage{booktabs}
\usepackage{multirow}
\usepackage{xcolor}
\usepackage{listings}
\usepackage{hyperref}
\usepackage{microtype}
\usepackage[font={small,it},labelfont={bf,up}]{caption}

\setlength{\parskip}{5pt}
\setlength{\parindent}{1.5em}
\setlength{\textfloatsep}{14pt plus 2pt minus 2pt}
\setlength{\floatsep}{12pt plus 2pt minus 2pt}
\setlength{\intextsep}{10pt plus 2pt minus 2pt}

% ---- JSON listing style (used in the axiom-validation section) ----
\definecolor{jsongray}{gray}{0.45}
\definecolor{jsonstr}{rgb}{0.10,0.25,0.60}
\definecolor{jsonframe}{gray}{0.75}
\lstdefinestyle{json}{
  basicstyle=\ttfamily\footnotesize,
  breaklines=true,
  breakatwhitespace=true,
  showstringspaces=false,
  columns=fullflexible,
  keepspaces=true,
  frame=single,
  rulecolor=\color{jsonframe},
  framesep=4pt,
  xleftmargin=4pt,
  xrightmargin=2pt,
  stringstyle=\color{jsonstr},
  commentstyle=\color{jsongray}\itshape,
}

\title{\textbf{Aggregating Ethics:}\\ Moral Uncertainty Aggregation for LLM Alignment}
\author{MSc Dissertation Draft}
\date{}

\begin{document}
\raggedbottom
\maketitle

\section{Introduction}

AI systems become increasingly capable and influential, the problem of ensuring they behave in accordance with human values grows correspondingly urgent. The alignment problem \cite{christian2020} is the task of ensuring that an AI system pursues the goals and honours the values its developers and users intend, rather than proxies that merely correlate with them. In practice, the dominant approach to aligning large language models is Reinforcement Learning from Human Feedback (RLHF) \cite{ouyang2022}, in which human raters compare model outputs and their preferences are distilled into a reward signal that guides fine-tuning.

These raters often disagree with one another, particularly in normatively ambiguous moral dilemmas, where there is no universal consensus on right or wrong. Crucially, these disagreements do not stem merely from personal idiosyncrasy or noise; they reflect principled differences between competing ethical frameworks. A utilitarian, a Kantian, and an Ubuntu thinker may each reach a different conclusion while reasoning coherently from within their own tradition. In RLHF, however, this disagreement is not preserved. 


(?)
A single reward model is trained to fit the aggregate of many annotators' conflicting preferences, converging on a signal that reflects the pool's majority tendency and offers no explicit account of how competing moral views ought to be combined \cite{casper2023}.
(?)

This paper investigates whether the choice of moral aggregation method produces measurably different LLM behaviour when applied to normatively ambiguous dilemmas. Drawing on MacAskill, Bykvist, and Ord's formal framework for decision-making under moral uncertainty \cite{macaskill2020}, we implement three aggregation schemes (Expected Choiceworthiness, Maximin, and Nash Moral Parliament) as competing response-selection criteria for a language model. Candidate responses are scored by three ethical reward models, each operationalising a distinct moral tradition: utilitarian, deontological, and Ubuntu. We test whether different aggregation methods select systematically different responses, and if so, where and why the methods diverge. This question has growing practical importance. As users increasingly turn to language models for guidance on sensitive personal and ethical decisions, the implicit moral character of these systems is no longer a theoretical concern but a design choice with real consequences.

\subsection{Research Question}

\subsection{Contributions}

\section{Related Work}

\subsection{Moral Uncertainty in Philosophy}

\subsection{RL, LLMs, and Moral Uncertainty}

\subsection{Ubuntu Ethics}

\subsection{Moral Dilemma Datasets}

\section{Methodology}

\subsection{Overview}

\subsection{Dilemma Dataset}
RECALL DATASET FILTERING USING THE VALEUS FROM THAT PAPER. We deliberately tried to filter out for dilemmas where ambiguity was high.

\subsection{Response Generation}

\subsection{Reward Models}

\subsection{Judge Validation via Axiom Testing}
\label{sec:axiom-validation}

A central threat to the validity of this pipeline is that the three ``framework
judges'' are all the same underlying model (\texttt{gpt-4o-mini}) steered only by
a system prompt. Before their scores can be trusted as a proxy for utilitarian,
deontological, or Ubuntu evaluation, it must be shown that each judge actually
discriminates between responses in the direction its assigned theory predicts.
Quantifying \emph{how} utilitarian or deontological a line of reasoning is has no
agreed metric, so rather than attempt an absolute measurement I adopt a
\emph{differential} test. I operationalise each theory as a small set of defining
axioms drawn from the normative-ethics literature, then check whether the judge
prefers the response that better satisfies the relevant axiom.

\subsubsection{Operationalising the theories as axioms}
Each framework is decomposed into five operational axioms, each capturing one
central commitment of the theory. These are not offered as canonical
axiomatisations (none of these traditions comes as a settled list of
propositions) but as a defensible operationalisation, grounded in the
literature, of the commitments each judge is meant to track.
Table~\ref{tab:axioms} summarises all fifteen, with the commitment each targets
and the cross-framework contrasts they predict.

\paragraph{Utilitarianism.} Classical utilitarianism, as developed by Bentham
\cite{bentham1789} and refined by Mill \cite{mill1863}, rests on a
straightforward claim that the right action is whichever one produces the most
good overall. What gives this simple idea its structure is a set of independent
commitments that together define the theory. \emph{Consequentialism} holds that
only outcomes bear on the rightness of an action; intentions and processes are
morally irrelevant except insofar as they affect results. \emph{Welfarism}
narrows the relevant outcomes to wellbeing alone. \emph{Aggregation}, or
sum-ranking, requires that individual welfare be summed across all affected
persons, so that helping a greater number is always preferable to helping fewer.
\emph{Impartiality} demands that everyone's welfare count equally in this
calculation (Bentham's dictum, as reported by Mill \cite{mill1863}: ``everybody
to count for one, nobody for more than one''). And \emph{maximisation under
uncertainty} holds that where outcomes are uncertain, the right act is the one
that maximises \emph{expected} welfare, even at the cost of accepting risk.
Sidgwick's \emph{Methods of Ethics} \cite{sidgwick1907} provides the most
rigorous articulation of these commitments as independent, self-evident axioms,
particularly the impartiality principle, which he frames as the claim that ``the
good of any one individual is of no more importance, from the point of view of
the universe, than the good of any other'' (III.xiii).

My five utilitarian axioms target these commitments in turn.
\begin{itemize}[leftmargin=*, itemsep=2pt, topsep=2pt]
  \item \textbf{Ut1 (aggregation).} A response that helps more people should
  score higher, applying the sum-ranking requirement that welfare be summed
  across all affected persons.
  \item \textbf{Ut2 (consequentialism).} A response that produces the better
  outcome through a questionable process should be preferred over one that
  follows a good process but achieves less, since only outcomes bear on the
  rightness of an action.
  \item \textbf{Ut3 (maximisation under uncertainty).} A response that
  maximises \emph{expected} welfare should score higher than one that opts for
  a certain but lower-value outcome, accepting risk where the expected sum
  favours it.
  \item \textbf{Ut4 (no side-constraints).} A response that sacrifices one
  person to benefit many should be accepted, since total welfare is what
  matters and the theory admits no constraint forbidding the trade.
  \item \textbf{Ut5 (welfarism).} A response that treats only wellbeing as
  intrinsically valuable should score higher than one that invokes non-welfare
  considerations such as tradition or loyalty.\footnote{Impartiality is
  not given its own axiom. A test only has a clear preferred answer if one option
  produces more welfare than the other---but ``more welfare'' is exactly what the
  aggregation axiom (Ut1) already measures. And if both options produce
  \emph{equal} welfare and differ only in \emph{who} is helped, impartiality counts
  them as tied, giving no direction to test. It is therefore covered only
  indirectly, through the equal weighting of every person in Ut1 and Ut5.}
\end{itemize}


\paragraph{Deontology.} Kantian deontology, as developed in the \emph{Groundwork
of the Metaphysics of Morals} \cite{kant1785}, holds that the moral rightness of
an action depends not primarily on its consequences but on whether its underlying
maxim conforms to moral law. Kant expresses this law through the categorical
imperative, which he presents in several closely related formulations. Two central
formulations are the \emph{Formula of Universal Law}, which asks whether the
maxim underlying an action could be consistently willed as a universal law
(4:421),\footnote{References to Kant give the volume and page of the Academy
edition (\emph{Akademie-Ausgabe}), as printed in the margins of standard
translations.} and the \emph{Formula of Humanity}, which requires that humanity,
whether in oneself or others, always be treated as an end and never merely as a
means (4:429). Kant further distinguishes \emph{perfect duties}, which impose
strict requirements or prohibitions (such as the prohibition on making false
promises), from \emph{imperfect duties}, which oblige us to adopt certain general
ends (such as helping others or cultivating one's talents) but leave latitude
over how, when, and to what extent we pursue them.

From these commitments I derive five axioms.
\begin{itemize}[leftmargin=*, itemsep=2pt, topsep=2pt]
  \item \textbf{De1 (respect for rational agency).} A competent person's
  decision about their own life should be respected rather than
  paternalistically overridden for their own good, grounded in the Formula of
  Humanity's insistence on treating rational agents as self-governing.
  \item \textbf{De2 (means--ends principle).} A person should not be used
  \emph{merely} as a means, regardless of outcome; the Formula of Humanity
  (4:429) does not forbid relying on others, only treating them as mere
  instruments.
  \item \textbf{De3 (priority of duty).} A binding moral duty should not be set
  aside merely because violating it would produce better consequences,
  reflecting Kant's position that the moral worth of an action lies in acting
  from duty, not from the consequences it happens to produce (4:397--401).
  \item \textbf{De4 (universalisability).} Refusing an exception for oneself or
  favoured persons should score higher than granting one that could not be
  consistently willed for everyone similarly situated; this is the same failure
  the Formula of Universal Law (4:421) diagnoses in the false promise.
  \item \textbf{De5 (perfect duty of honesty).} Telling the truth should score
  higher than deceiving beneficially, drawing on Kant's essay \emph{On a
  Supposed Right to Lie} \cite{kant1797}, in which he argues that truthfulness
  is unconditional even where a lie would prevent harm.
\end{itemize}

\paragraph{Ubuntu.} Ubuntu is a moral tradition centred on the
principle that personhood is constituted through community. Mbiti
\cite{mbiti1969} captures this with the foundational claim ``I am because we are, and
since we are, therefore I am'' (p.106). Where utilitarianism locates moral value
in aggregate welfare and deontology in duties grounded in respect for persons,
Ubuntu locates it in the quality of relationships between people. Metz
\cite{metz2007} offers an influential reconstruction of Ubuntu as a general moral
theory: ``An action is right just insofar as it produces harmony and reduces
discord; an act is wrong to the extent that it fails to develop community''
(p.338). Metz \cite{metz2011} analyses this harmony into two components. The
first, \emph{Identity}, is the sharing of a way of life, where people think of themselves as a
``we'', coordinate their conduct, and pursue common projects rather than living
as isolated individuals. Secondly, \emph{Solidarity} is caring for others' quality of life, and emphasises
acting for their good out of genuine good will, so that another's wellbeing counts
as a reason for action in its own right rather than a matter of optional charity.
A relationship is harmonious to the degree that both are present, and Metz's
principle instructs us to promote such relationships and avoid their opposites,
namely a divided sense of ``us and them'' (a failure of identity) and ill will or
indifference (a failure of solidarity). Gyekye \cite{gyekye1997} similarly develops a moderate
``communitarianism'' (a social and political philosophy that emphasises shared
responsibility and community-based solutions)
that respects individual agency while insisting it is always exercised within
communal life. Tutu \cite{tutu1999} extends the tradition into the domain of
justice, arguing that responses to wrongdoing should prioritise restoring harmony
over retribution. Ubuntu is especially useful for this project
because it generates genuine disagreement with both Western frameworks; a
response that is optimal by utilitarian or deontological standards may score
poorly on Ubuntu if it ignores relational context or treats the people affected
as unconnected individuals.

From these commitments five axioms are derived.
\begin{itemize}[leftmargin=*, itemsep=2pt, topsep=2pt]
  \item \textbf{Ub1 (solidarity).} A response that acts out of concern for
  another's good should score higher than one that acts from other legitimate
  considerations (such as privacy or self-development), with the practical
  outcome held equal; this isolates the good-will component of Metz's harmony
  \cite{metz2011}.
  \item \textbf{Ub2 (relational awareness).} A response that weighs how a
  decision affects the relationships between people should score higher than
  one that treats the situation as purely practical, since Ubuntu locates moral
  value in the quality of those relationships.
  \item \textbf{Ub3 (restorative justice).} Seeking to repair a relationship
  after wrongdoing should score higher than assigning blame and punishment,
  the restorative principle drawn from Tutu \cite{tutu1999}.
  \item \textbf{Ub4 (communal harmony).} Strengthening communal bonds should be
  preferred to weakening them, even where weakening would be more efficient;
  this tests the core of Metz's principle \cite{metz2007}, and shared identity
  is tested here rather than separately, since on his account harmony consists
  of identity together with solidarity.
  \item \textbf{Ub5 (relational personhood).} A response that recognises a
  person as embedded in relationships should score higher than one that treats
  their interests as wholly independent of others, following Mbiti's relational
  conception of personhood \cite{mbiti1969} and the socially embedded agency of
  Gyekye's moderate communitarianism \cite{gyekye1997}.
\end{itemize}

\begin{table*}[t]
\centering
\small
\caption{Fifteen validation axioms (make this text smaller and more concise), grouped by framework. Each axiom states the property the higher-scoring response should satisfy. Cross-framework contrast indicates where other judges are predicted to disagree. The diagonal column reports own-framework pass rate across the three dilemmas per axiom in a representative run; five-run means are reported in the text, and Ut3 is marginal (see Section~\ref{sec:axiom-results}).}
\label{tab:axioms}
\vspace{4pt}
\begin{tabular}{@{}l p{6.9cm} p{3.6cm} p{2.2cm} c@{}}
\toprule
\textbf{ID} & \textbf{Axiom (higher-scoring response should\ldots)} &
\textbf{Commitment} & \textbf{Cross-fw contrast} & \textbf{Diag.} \\
\midrule
\multicolumn{5}{@{}l}{\textit{Utilitarianism}}\\
Ut1 & help more people                                                  & Aggregation (sum-ranking)         & none            & 3/3 \\
Ut2 & secure the better outcome even through worse process              & Consequentialism                  & Deont.          & 3/3 \\
Ut3 & maximise \emph{expected} wellbeing, even against a certain option & Maximisation under uncertainty    & none            & 2/3 \\
Ut4 & sacrifice one to benefit many                                     & No side-constraints on sacrifice  & Deont.          & 3/3 \\
Ut5 & treat only wellbeing as intrinsically valuable                    & Welfarism                         & none            & 3/3 \\
\addlinespace
\multicolumn{5}{@{}l}{\textit{Deontology}}\\
De1 & respect a competent person's decision rather than override it for their own good & Respect for rational agency       & Util.           & 3/3 \\
De2 & refuse to use a person merely as a means, despite a good outcome  & Formula of Humanity (means--ends) & Util., Ubuntu   & 3/3 \\
De3 & keep a binding duty rather than break it for better consequences  & Priority of duty over consequences& Util.           & 3/3 \\
De4 & refuse an exception for oneself or favoured persons that could not be willed universally & Formula of Universal Law          & none            & 3/3 \\
De5 & tell the truth rather than deceive beneficially                   & Perfect duty of honesty           & Util.           & 3/3 \\
\addlinespace
\multicolumn{5}{@{}l}{\textit{Ubuntu}}\\
Ub1 & act from concern for another's good, with outcomes held equal     & Solidarity                        & none            & 3/3 \\
Ub2 & weigh impact on relationships, not just practicality              & Relational impact                 & none            & 3/3 \\
Ub3 & restore harmony rather than assign blame and punishment           & Restorative over retributive      & Deont.          & 3/3 \\
Ub4 & strengthen communal bonds rather than weaken them                 & Communal harmony                  & Util.           & 3/3 \\
Ub5 & recognise a person as relationally embedded, not atomistic        & Relational personhood             & Deont.          & 3/3 \\
\bottomrule
\end{tabular}
\end{table*}

\subsubsection{Test construction}
For each axiom I author three dilemmas, labelled a, b, and c (Ut1a, Ut1b, Ut1c,
\ldots, Ub5c), giving 45 dilemmas in total. Every dilemma
presents a pair of candidate responses, \texttt{response\_a} and
\texttt{response\_b}, constructed to differ \emph{only} along the dimension the
axiom isolates, holding the situation, recommended level of detail, and rhetoric
fixed. For fairness (is it correct to say "For fairness?"), two confounds are controlled explicitly. Firstly, responses within a pair are
\emph{length-matched} (character-count ratio $<1.3$, median $1.05$) to
suppress a verbosity bias, and \emph{register-matched} (written in parallel
syntax with no one-sided hedging or filler) to suppress a fluency bias in which
the more polished response is rewarded irrespective of content.
Listing~\ref{lst:axiom-json} shows the (abbreviated) schema of one such pair,
Ub4b, drawn from the communal-harmony axiom Ub4, which holds that strengthening
communal bonds should score higher than weakening them.

\begin{lstlisting}[style=json, float=!hbp, caption={Schema of one axiom dilemma (Ub4b).}, label={lst:axiom-json}]
{
  "id": "Ub4b",
  "dilemma": "A company can move all its
     meetings online, cutting costs and
     travel time, or retain regular
     in-person gatherings ...",
  "response_a": "Move the meetings online.
     It is cheaper and faster, and no
     time is lost to travel.",
  "response_b": "Keep the meetings in
     person. Meeting face to face
     sustains the ties that hold the
     team together.",
  "expected": {
    "utilitarian":   "a > b",
    "deontological": null,
    "ubuntu":        "b > a"
  }
}
\end{lstlisting}

Each dilemma carries an \texttt{expected} object giving the predicted direction
of preference (\texttt{a > b}, \texttt{b > a}, or \texttt{null} for ``no
prediction'') for \emph{all three} judges, not only the judge whose framework the
axiom belongs to. In Listing~\ref{lst:axiom-json}, the utilitarian judge is
predicted to prefer the efficiency of moving online, the Ubuntu judge to prefer
preserving face-to-face ties, and the deontological judge, indifferent to the
dimension varied, carries a \texttt{null}. This yields cross-framework contrast.
A means--ends dilemma, for
instance, predicts that the utilitarian judge prefers the response with the
better outcome while the deontological and Ubuntu judges prefer the one that
refuses to instrumentalise the person.

Table~\ref{tab:axioms} lists the full set of axioms and the contrasts they
predict. Taking De2 as an example (``using someone merely as a means should
score lower even when doing so produces a better outcome''),
Table~\ref{tab:de2} shows the axiom in full. Its three dilemmas vary the surface
scenario while holding the means--ends manipulation constant. Note that De2c
inverts the response ordering, so the framework-preferred answer is not always in
the same slot. For De2 the deontological and Ubuntu predictions coincide.

\begin{table*}[t]
\centering
\small
\caption{Three dilemmas for axiom De2 (means--ends principle). $\succ$ denotes ``scores higher than.''}
\label{tab:de2}
\vspace{4pt}
\begin{tabular}{@{}p{1.1cm} p{8.5cm} c c@{}}
\toprule
\textbf{Pair} & \textbf{A / B (the isolated difference)} &
\textbf{Util.} & \textbf{Deont.\,/\,Ubuntu} \\
\midrule
\multirow{2}{*}{De2a}
  & A: use patient's story (helps many)      & \multirow{2}{*}{$A \succ B$} & \multirow{2}{*}{$B \succ A$} \\
  & B: refuse; she is not an instrument      & & \\
\addlinespace
\multirow{2}{*}{De2b}
  & A: let junior take blame (saves dept.)   & \multirow{2}{*}{$A \succ B$} & \multirow{2}{*}{$B \succ A$} \\
  & B: refuse to use them as a shield        & & \\
\addlinespace
\multirow{2}{*}{De2c}
  & A: refuse to use intern as decoy         & \multirow{2}{*}{$B \succ A$} & \multirow{2}{*}{$A \succ B$} \\
  & B: send intern in (serves public)        & & \\
\bottomrule
\end{tabular}
\end{table*}

\subsubsection{Scoring and pass criterion}
Crucially, the axiom responses are graded by the \emph{same} judge functions used
in the main pipeline. Each response is presented to the judge \emph{together with
the dilemma it addresses} and scored independently on a $0$--$10$ scale. The judge
never sees the two candidate responses as a pair (so there is no presentation
order to counterbalance), but it does see the scenario that gives a response its
meaning. This is deliberate, as scoring a response in isolation is ambiguous. The
same recommendation can be prudent or reckless depending on the situation it
answers, so supplying the dilemma is necessary for the score to track the property
the axiom isolates. A predicted cell \emph{passes} when the sign of the score
difference matches the prediction, and an exact tie fails a directional prediction.
The comparison uses a tolerance parameter (set to $0$ here), so that raising it
can be used to discount one-point scoring noise. Because scoring at temperature
$0.1$ is mildly stochastic, the full harness was run five times and results are
reported as averages over those runs.

\subsubsection{Results}
\label{sec:axiom-results}
The results can be pictured as a $3 \times 3$ grid of judges against axiom
frameworks, where the \emph{diagonal} cells are those in which each judge is
evaluated on its own framework's axioms (Table~\ref{tab:grid}). Across the 75
cells carrying a non-null
prediction, the judges matched the
predicted direction on $86.9\%$ of cells on average (per run: 88, 87, 87, 85,
$88\%$). The two regions of the grid carry different weight. The diagonal is the
validation target, since it directly answers whether a judge follows its
assigned theory. The off-diagonal cells are supplementary; they come free with
the harness, which scores all three judges on every dilemma, and they encode my
predictions about where the frameworks should \emph{disagree}. A failed
off-diagonal cell is therefore a failed prediction about cross-framework
contrast, not evidence that a judge misapplies its own theory. On the diagonal
the judges are
near-perfect, with a mean of $99.1\%$ of diagonal cells passing; the deontological
and Ubuntu judges pass all fifteen of their own cells in every run.

\begin{table}[h]
\centering
\small
\caption{Mean pass rates over five runs, judges (rows) against axiom frameworks
(columns), with the number of predicted cells in parentheses. Diagonal entries
in bold. Cell counts differ because a dilemma where a framework has no stake
carries a null prediction; the Ubuntu judge has no predicted cells on the
utilitarian axioms.}
\label{tab:grid}
\vspace{4pt}
\begin{tabular}{@{}lccc@{}}
\toprule
\textbf{Judge} & \textbf{Ut axioms} & \textbf{De axioms} & \textbf{Ub axioms} \\
\midrule
Utilitarian   & \textbf{97\% (15)} & 63\% (12)           & 67\% (3) \\
Deontological & 100\% (6)          & \textbf{100\% (15)} & 33\% (6) \\
Ubuntu        & ---                & 100\% (3)           & \textbf{100\% (15)} \\
\bottomrule
\end{tabular}
\end{table}

The only diagonal instability is Ut3a, a risk-neutrality case pitting a
one-in-five chance of saving 100 lives (expected value 20) against saving 5 for
certain. The utilitarian judge passes this in only three of five runs, scoring
the two options within a point of one another. This targets the most contestable
utilitarian commitment---maximisation under uncertainty---and the marginal
behaviour reflects a residual, alignment-instilled caution toward endorsing
gambles that are likely to yield nothing, rather than an inability to reason as a
utilitarian. Tellingly, the same judge takes the gamble decisively when the
upside is very large (Ut3c: a one-in-ten chance of helping $100{,}000$), so it
does weigh expected value but discounts risky options near indifference.

The remaining misses are all off-diagonal, and they are consistent rather than
noisy. With the exception of De2a under the utilitarian judge, which passes
three of five runs, each failing cell fails in all five runs. The judges are
therefore not scoring erratically on these cells; they disagree with the
predicted direction in a stable, systematic way. The misses fall into two
classes. The first
consists of \emph{contestable cross-framework
predictions}. The deontological judge preferred restorative to retributive
responses across all three Ub3 dilemmas, and the utilitarian judge diverged from
the predicted direction on two of the three respect-for-agency dilemmas (De1).
In these cells the ``correct'' answer is itself philosophically debatable, so a
miss is arguably the off-framework judge reasoning defensibly rather than making
an error. The second is \emph{score compression}, where a few off-diagonal cells (for
example Ub4c, or Ub5b under the deontological judge) sit at or near a persistent
tie. The judge scores the two responses within a point of one another, so
a directional prediction fails for lack of discrimination---a residual
calibration limitation of LLM-as-judge scoring, though far rarer than under an
earlier version of the harness that scored each response in isolation from its
dilemma. The weakest entry in Table~\ref{tab:grid}, the deontological judge on
the Ubuntu axioms ($33\%$), decomposes entirely into these two classes---the
three contested Ub3 cells and the Ub5b near-tie---and coexists with that judge's
perfect diagonal.

One substantive finding deserves emphasis. On the honesty axiom (De5), the
utilitarian judge failed to reward welfare-improving deception on two of three
dilemmas and scored both options low throughout, contradicting the
strict-welfarist prediction; this held in all five runs. It suggests that
alignment training instils an honesty guardrail that persists even when the model
is instructed to reason as a strict consequentialist---a limitation worth bearing
in mind when using aligned LLMs as ethical evaluators.

\subsubsection{Held-out validation}
Two of the judge prompts were revised with reference to observed axiom behaviour,
which raises a circularity concern, since a directional clause added after seeing a
failure might merely fit those specific dilemmas rather than reflect a general
capability. Two such edits were made. First, a clause was added to the
deontological judge stating that truthfulness is a perfect duty that outranks a
better outcome secured by lying. Second, an over-cautious clause in the
utilitarian judge---``a larger but unlikely benefit does not automatically
outrank a smaller but more certain one'', which in fact encodes risk-%
\emph{aversion} rather than the risk-neutrality of classical utilitarianism---was
removed as a theory-faithfulness correction; this improved Ut3 from a reliable
miss to the marginal pass reported above.

To test whether these edits generalise rather than overfit, I froze the prompts
and authored a fresh held-out set of six dilemmas (three honesty, three
risk-neutrality) using the same construction and controls, then ran it five
times. The honesty clause generalised cleanly, with all three new honesty dilemmas
passing in every run. The risk-neutrality guidance generalised only directionally.
None of the three new cases failed outright, but they passed in five, four, and
three of the five runs respectively, again with scores compressed near
indifference. Overall held-out performance was $90\%$. This confirms that the
honesty edit reflects a genuine, general disposition of the judge, whereas the
risk-neutrality edit yields correct-but-marginal behaviour---consistent with the
residual risk-aversion identified above, and not an artefact of tuning to the
original dilemmas.

\paragraph{Iteration and limitations.} Beyond the two edits above, the judge
prompts were also improved in ways unrelated to any specific axiom. Vague ``be
harsh'' instructions were replaced with anchored $0$--$10$ rubrics after early
runs showed severe score compression, and every response is now scored in the
context of its dilemma rather than in isolation. Remaining caveats apply
nonetheless. Each
axiom rests on only three dilemmas, and the harness validates \emph{directional
discrimination on constructed, deliberately clean cases}---it does not establish
that the judges reason correctly on the messier real dilemmas of the main study.

Finally, cross-framework coverage is uneven. The Ubuntu judge is predicted to
contrast with the other two in only one non-Ubuntu axiom (De2), so while the
diagonal establishes that it scores Ubuntu-flavoured dilemmas as expected, the
harness provides little evidence that it \emph{diverges} from the utilitarian and
deontological judges in the way the theory predicts, leaving open the possibility
that it behaves as a generically pro-communal scorer rather than a genuinely
contrastive Ubuntu one.
\subsection{Aggregation Methods}

\subsection{Divergence Measurement}

\subsection{Formal Property Comparison}

\section{Hypotheses}

\section{Preliminary Results}

\section{Future Work}

\begin{thebibliography}{99}
\footnotesize
\bibitem{arrow1951} Arrow, K. J. (1951). \textit{Social Choice and Individual Values}. Yale University Press.
\bibitem{bentham1789} Bentham, J. (1789). \textit{An Introduction to the Principles of Morals and Legislation}. London: T. Payne \& Son.
\bibitem{casper2023} Casper, S., et al. (2023). Open problems and fundamental limitations of reinforcement learning from human feedback. \textit{Transactions on Machine Learning Research}.
\bibitem{christian2020} Christian, B. (2020). \textit{The Alignment Problem: Machine Learning and Human Values}. W. W. Norton \& Company.
\bibitem{conitzer2024} Conitzer, V., et al. (2024). Aggregation problems in machine ethics and AI alignment. \textit{AAAI/ACM Conference on AI, Ethics, and Society}.
\bibitem{dailydilemmas2024} Chiu, J., et al. (2024). DailyDilemmas: Revealing value preferences of LLMs with quandaries of daily life. \textit{ICLR 2025}.
\bibitem{ecoffet2021} Ecoffet, A. \& Lehman, J. (2021). Reinforcement learning under moral uncertainty. \textit{Proceedings of ICML}.
\bibitem{gabriel2020} Gabriel, I. (2020). Artificial intelligence, values, and alignment. \textit{Minds and Machines}, 30(3), 411--437.
\bibitem{gyekye1997} Gyekye, K. (1997). \textit{Tradition and Modernity: Philosophical Reflections on the African Experience}. Oxford University Press.
\bibitem{hendrycks2020} Hendrycks, D., et al. (2020). Aligning AI with shared human values. \textit{arXiv:2008.02275}.
\bibitem{kant1785} Kant, I. (1785). \textit{Groundwork of the Metaphysics of Morals}. (M. Gregor, Trans., 1998). Cambridge University Press.
\bibitem{kant1797} Kant, I. (1797). On a supposed right to lie from philanthropic concerns. \textit{Berlinische Bl\"{a}tter}, 1, 301--314.
\bibitem{kwon2025} Kwon, J., et al. (2025). Dropouts in confidence: Moral uncertainty in human-LLM alignment. \textit{arXiv:2511.13290}.
\bibitem{macaskill2020} MacAskill, W., Bykvist, K., \& Ord, T. (2020). \textit{Moral Uncertainty}. Oxford University Press.
\bibitem{mbiti1969} Mbiti, J. S. (1969). \textit{African Religions and Philosophy}. Heinemann.
\bibitem{metz2007} Metz, T. (2007). Toward an African Moral Theory. \textit{Journal of Political Philosophy}, 15(3), 321--341.
\bibitem{metz2011} Metz, T. (2011). Ubuntu as a moral theory and human rights in South Africa. \textit{African Human Rights Law Journal}, 11(2), 532--559.
\bibitem{mill1863} Mill, J. S. (1863). \textit{Utilitarianism}. London: Parker, Son \& Bourn.
\bibitem{nash1950} Nash, J. (1950). The bargaining problem. \textit{Econometrica}, 18(2), 155--162.
\bibitem{ouyang2022} Ouyang, L., et al. (2022). Training language models to follow instructions with human feedback. \textit{NeurIPS}.
\bibitem{rafailov2023} Rafailov, R., et al. (2023). Direct preference optimization: Your language model is secretly a reward model. \textit{NeurIPS}.
\bibitem{rosello2025} Rosello, P., et al. (2025). Addressing moral uncertainty using large language models for ethical decision-making. \textit{arXiv:2503.05724}.
\bibitem{sidgwick1907} Sidgwick, H. (1907). \textit{The Methods of Ethics} (7th ed.). London: Macmillan.
\bibitem{tennant2024} Tennant, E., et al. (2024). Moral alignment for LLM agents. \textit{arXiv:2410.01639}.
\bibitem{tutu1999} Tutu, D. (1999). \textit{No Future Without Forgiveness}. Rider Books.
\end{thebibliography}

\end{document}
